%%%%%%%%%%%%%%%%%%%%%%%%%%%%%%%%%%%%%%%%%%%%%%%%%%%%%%%%%%%%%%%%%%%%%%%%%%%%%%%
%% scope_box.tex -- the boxed scope statement. Input TWICE, deliberately:
%%   (1) Introduction, after the contributions/regime table, so a referee meets the
%%       scope before the results;
%%   (2) \S Limitations, as the opening block.
%% It contains no \label, so a second \input is safe.
%%
%% PORTABILITY: \fbox+minipage rather than tcolorbox, so the same source compiles in the
%% IEEEtran two-column builds (kbound.tex, kbound_short.tex) and the single-column TMLR
%% build (kbound_tmlr.tex). Requires nothing beyond base LaTeX.
%%%%%%%%%%%%%%%%%%%%%%%%%%%%%%%%%%%%%%%%%%%%%%%%%%%%%%%%%%%%%%%%%%%%%%%%%%%%%%%

\begin{center}
\fbox{%
\begin{minipage}{0.94\columnwidth}
\small
\textbf{Scope of the guarantee.} K-Bound does not guarantee that a newly encountered
deployment is exchangeable with the calibration data, that label-free evidence is
risk-aligned, or that a fitted benefit estimator transfers. Its guarantee is conditional:
\emph{if} the declared benefit interval attains marginal target coverage, \emph{then}
strict adaptation on $\widehat\Delta-\varepsilon>0$ controls the unconditional
false-adapt event at level $\alpha$, and strict freezing on
$\widehat\Delta+\varepsilon<0$ the false-freeze event. Deployment diagnostics can
identify support shift, instability, leakage and insufficient effective sample size, but
passing them does not prove coverage. When required checks fail or cannot be evaluated,
the prescribed output is \textsc{abstain} or \textsc{freeze} --- not a certified
adaptation decision.
\end{minipage}}
\end{center}
